% ICRA 2027 double-anonymous draft. All quantitative claims are traceable to
% code, checkpoints, or recorded demonstrations in the ZQ branch.
\documentclass[letterpaper,10pt,conference]{ieeeconf}
\IEEEoverridecommandlockouts
\overrideIEEEmargins

\usepackage{amsmath,amssymb}
\usepackage{graphicx}
\usepackage{booktabs}
\usepackage{multirow}
\usepackage{xcolor}
\usepackage{cite}
\usepackage{url}

\newcommand{\TODO}[1]{\textcolor{red}{\textbf{[TODO: #1]}}}
\newcommand{\obs}{\mathbf{o}}
\newcommand{\act}{\mathbf{a}}
\newcommand{\goal}{g}
\newcommand{\hid}{\mathbf{h}}

\title{MOSAIC-Broncho: Multimodal Observation and State-Aware\\
Imitation Control for Goal-Directed Autonomous Bronchoscopy}

\author{Anonymous Authors}

\begin{document}
\maketitle
\thispagestyle{empty}
\pagestyle{empty}

\begin{abstract}
A bronchial bifurcation often looks clear to a human but remains ambiguous to a
robot. The image alone does not reveal the intended branch, recent motion, or
tendon hysteresis. We therefore present MOSAIC-Broncho, a Multimodal
Observation and State-Aware Imitation Control framework for goal-directed
bronchoscopy. A unified monocular front end extracts structural, anatomical,
and metric-geometric cues. These cues are synchronized with motor position,
finite-difference velocity, and frame timing to form an auditable 20-D state.
The controller conditions an eight-frame history on the desired anatomical
target and predicts bounded increments for insertion and two bending axes. It
is trained from 87 real-robot demonstrations with 30,730 synchronized frames
across seven targets. The resulting 773,221-parameter GRU reaches a held-out
action MSE of 0.002084 at epoch 29. The study evaluates each link in the loop:
perception quality, temporal and multimodal ablations, runtime behavior, and
target-wise navigation in a silicone airway phantom. The design favors a small,
inspectable state and continuous real-time control over a large end-to-end
policy. It provides a practical route toward autonomous bronchoscopy when data,
compute, and safety margins are limited.
\end{abstract}

\begin{keywords}
Medical robots and systems; imitation learning; recurrent neural networks;
autonomous navigation; bronchoscopy; multimodal perception.
\end{keywords}

\section{Introduction}

A bronchoscope reaches a fork. Both openings may look plausible. Yet only one
leads to the target. The controller must remember the route, read the local
geometry, and account for what the instrument has just done. A single image is
not enough.

This compact example captures the central difficulty of autonomous
bronchoscopy. Airway views are repetitive. Their appearance changes with
reflection, mucus, deformation, and blur. The robot adds another source of
ambiguity: tendon friction and hysteresis make the same command produce
different motion histories \cite{webster2010continuum,camarillo2008mechanics}.
Although robotic platforms improve reach and stability, operation remains
largely expert driven \cite{fielding2019firsthuman,chaddha2019multicenter,
rojassolano2018feasibility}.

Recent work attacks different parts of this problem. Visual servoing keeps the
lumen near the image center \cite{lazo2022visualservo}. Online topological
localization estimates where the scope is without patient-specific CT
\cite{tomasini2025topological}. AI co-pilot learns shared control from a virtual
expert and uses DAgger to reduce distribution shift
\cite{zhang2024aicopilot}. BronchoCopilot fuses images and estimated pose in a
reinforcement-learning agent \cite{zhao2024bronchocopilot}. More recently,
long--short agents have combined a fast visual controller with slower strategic
reasoning, reporting phantom, ex-vivo, and in-vivo evaluation
\cite{wu2026longshort}. Together, these systems show rapid progress. They also
leave a practical gap. Simulation offers scale but creates transfer risk.
Purely reactive control loses route context. Large or hierarchical policies can
be difficult to run and audit at motor-loop rate.

Our premise is simple. Do not ask one image to carry the whole navigation
problem. Instead, expose the few signals that matter and let a recurrent policy
reason over their recent history. MOSAIC-Broncho follows this premise. Its
MonoAirway-Fusion (MAF) front end turns a monocular frame into structural,
anatomical, and metric-geometric evidence. The policy then combines these cues
with the commanded goal and the measured motor state. It predicts continuous
increments for insertion and two bending axes.

This is not a claim that a GRU is novel by itself. GRUs are established sequence
models \cite{cho2014gru}, and goal conditioning is well studied
\cite{codevilla2018conditional}. The contribution lies in the control
formulation: a small, normalized, and auditable state connects monocular airway
perception to real three-axis demonstrations. This design is deliberately
data-efficient. It is also honest about the limits of behavior cloning, notably
covariate shift outside expert-supported states
\cite{argall2009survey,osa2018imitation,ross2011dagger}.

We ask whether this compact formulation is sufficient for goal-directed
navigation on a physical bronchial phantom. Our contributions are:
\begin{enumerate}
  \item a unified monocular perception interface that exposes complementary
  structural, anatomical, and metric-geometric evidence;
  \item a normalized 20-D state that joins this evidence with motor position,
  velocity, and timing, making every policy input loggable and inspectable;
  \item a 773k-parameter, goal-conditioned GRU that maps eight-frame histories
  from real expert demonstrations to bounded three-axis commands; and
  \item an evaluation design that connects component quality and controlled
  ablations to target-wise navigation on a physical silicone phantom.
\end{enumerate}

\section{Results}

\subsection{Unified Multimodal Perception}

The structural output of MAF was evaluated for bronchial-junction segmentation in silicone
airway images. It correctly delineated bifurcation regions in representative
clean, low-light, specular, and blurred views. Quantitative reporting will use
a fixed held-out test split and include Dice, IoU, precision, recall, and
inference latency; these values are reserved for the recovered evaluation
sheet rather than estimated from qualitative examples.

The nine-class anatomical output head contains 2.59M parameters. Its saved
validation results are precision 0.8926, recall 0.7797, mAP50 0.9096, and
mAP50--95 0.6006. The gap between precision and recall motivates the temporal
confirmation logic used by the runtime anatomical state estimator: isolated
missed frames do not immediately invalidate the confirmed airway location.

\begin{table}[t]
\caption{MAF Output Quality in Silicone Airways}
\label{tab:perception-results}
\centering\small
\resizebox{\columnwidth}{!}{%
\begin{tabular}{lccccc}
\toprule
Output & Precision & Recall & Dice/mAP50 & IoU/mAP50--95 & Latency \\
\midrule
Structure & \TODO{--} & \TODO{--} & \TODO{Dice} & \TODO{IoU} & \TODO{ms} \\
Anatomy & 0.8926 & 0.7797 & 0.9096 & 0.6006 & \TODO{ms} \\
\bottomrule
\end{tabular}}
\end{table}

\subsection{Metric-Geometric Perception}

The metric-geometric output of MAF was tested inside the laboratory silicone bronchial phantom
at predefined camera-to-surface distances and at multiple airway locations.
We report both raw model output and the linearly calibrated estimate. The final
analysis will include sample count, mean absolute error (MAE), root mean square
error (RMSE), signed bias, coefficient of determination, and valid-frame rate.
Performance will also be stratified by normal illumination, specular
reflection, partial occlusion, and motion blur.

\begin{table}[t]
\caption{MAF Metric-Geometric Evaluation}
\label{tab:depth-results}
\centering\small
\resizebox{\columnwidth}{!}{%
\begin{tabular}{lcccc}
\toprule
Output & $N$ & MAE (mm) & RMSE (mm) & Bias (mm) \\
\midrule
Raw & \TODO{--} & \TODO{--} & \TODO{--} & \TODO{--} \\
Calibrated & \TODO{--} & \TODO{--} & \TODO{--} & \TODO{--} \\
\bottomrule
\end{tabular}
}
\end{table}

\subsection{Recurrent Policy Learning and Ablation}

The real-robot dataset contains 87 expert demonstrations and 30,730
synchronized frames across seven recorded anatomical targets. The selected
MOSAIC-GRU checkpoint contains 773,221 trainable parameters. It was
selected at epoch 29 with a held-out action MSE of 0.0020837828 and auxiliary
task accuracy of 1.000. The first logged validation action MSE was 0.02577,
corresponding to a 91.9\% reduction at the selected checkpoint. Task accuracy
must not be interpreted as navigation success because it measures only the
recorded navigate/intervene label.

The ablation study evaluates the full model against a frame-wise MLP, a GRU
without goal conditioning, a single-head GRU, a model without depth/path
features, and a model without velocity/timing features. Every variant uses the
same trajectory-wise split and reports total and per-axis action error,
macro-F1, parameter count, and hardware route success.

\begin{table}[t]
\caption{MOSAIC-GRU Ablation}
\label{tab:gru-ablation}
\centering\small
\resizebox{\columnwidth}{!}{%
\begin{tabular}{lcccc}
\toprule
Variant & Action MSE & Task F1 & Parameters & Route success \\
\midrule
Frame-wise MLP & \TODO{--} & \TODO{--} & \TODO{--} & \TODO{--} \\
GRU, no goal & \TODO{--} & \TODO{--} & \TODO{--} & \TODO{--} \\
GRU, action only & \TODO{--} & -- & \TODO{--} & \TODO{--} \\
No depth/path & \TODO{--} & \TODO{--} & \TODO{--} & \TODO{--} \\
No velocity/$\Delta t$ & \TODO{--} & \TODO{--} & \TODO{--} & \TODO{--} \\
Full model & 0.002084 & \TODO{macro-F1} & 773,221 & \TODO{--} \\
\bottomrule
\end{tabular}}
\end{table}

\subsection{Real-Robot Navigation Success}

The complete policy was evaluated from a fixed start pose in the silicone
bronchial phantom. A trial is successful only if the bronchoscope enters the
predefined target branch without operator takeover and remains there for the
confirmation interval. We report per-target and pooled success with Wilson
95\% confidence intervals, together with wrong-branch rate, completion time,
contact count, and intervention count.

\begin{table}[t]
\caption{Real-Robot Navigation Results}
\label{tab:evidence}
\centering\small
\resizebox{\columnwidth}{!}{%
\begin{tabular}{lll}
\toprule
Target & Success ($n/N$) & 95\% CI \\
\midrule
LMB & \TODO{--/--} & \TODO{--} \\
RMB & \TODO{--/--} & \TODO{--} \\
LUB & \TODO{--/--} & \TODO{--} \\
LLB & \TODO{--/--} & \TODO{--} \\
RUL & \TODO{--/--} & \TODO{--} \\
RML & \TODO{--/--} & \TODO{--} \\
RLL & \TODO{--/--} & \TODO{--} \\
Pooled & \TODO{--/--} & \TODO{--} \\
\bottomrule
\end{tabular}}
\end{table}

Together, these four result groups evaluate the complete perception-to-action
chain: unified multimodal perception, recurrent policy learning, controlled
ablation, and physical target arrival. Values shown as placeholders require
the corresponding saved evaluation sheet and are not inferred from videos.

\section{System and Multimodal Observation}

\subsection{Robotic Platform}

The platform uses a flexible bronchoscope driven by three position-controlled
motors: insertion/withdrawal $M_0$, left--right bending $M_1$, and up--down
bending $M_2$. The software separates camera/perception and motor-control
threads through a thread-safe feature dictionary. Commands are checked against
the implemented limits $M_0\in[-900,0]^\circ$, $M_1\in[-170,170]^\circ$, and
$M_2\in[-600,600]^\circ$. The policy interface accepts the newest perception
snapshot and cached motor feedback and can be called at 20--60\,Hz.

\subsection{MonoAirway-Fusion Perception}

MAF treats monocular perception as one structured inference operator. For an
RGB frame $I_t$, it exposes structural evidence $\mathbf s_t$, anatomical
evidence $\mathbf c_t$, and metric-geometric evidence $\mathbf d_t$:
\begin{equation}
 (\mathbf s_t,\mathbf c_t,\mathbf d_t)=\Phi_{\rm MAF}(I_t).
\end{equation}
The structural stream provides obstruction and bifurcation centroids, areas,
and validity; the anatomical stream provides airway class and confidence; the
metric-geometric stream provides center, mean, and maximum depth together with
left, center, and right free-space summaries. The streams form a single
13-dimensional visual-geometric interface rather than high-dimensional maps.

For reproducibility, the present implementation instantiates the structural
stream with a VGG--U-Net segmentation backbone \cite{ronneberger2015unet}, the
anatomical stream with a lightweight YOLO-family detector
\cite{redmon2016yolo}, and the monocular-depth prior with Depth Anything V2
\cite{yang2024depthanythingv2}. These are replaceable estimators inside MAF;
the method-level contribution is their synchronized, physically interpretable
interface and its coupling to recurrent control.

Let $\mathbf v_t=[\mathbf s_t;\mathbf c_t;\mathbf d_t]\in\mathbb R^{13}$.
For each time $t$, $\obs_t\in\mathbb{R}^{20}$ contains seven visual features,
six depth/path features, three motor angles, three finite-difference motor
velocities, and the elapsed time $\Delta t$. Each component is normalized by a
fixed physical or geometric scale $r_i$:
\begin{equation}
 \tilde o_{t,i}=o_{t,i}/r_i,\qquad
 \tilde\obs_t=[\mathbf v_t;\mathbf q_t;\dot{\mathbf q}_t;\Delta t].
\end{equation}
The resulting representation is small
enough to inspect, log, and replay, while temporal ambiguity is delegated to
the recurrent policy.

\begin{table}[t]
\caption{Structured Observation Used by the Deployed Policy}
\label{tab:obs}
\centering\small
\begin{tabular}{lcl}
\toprule
Group & Dim. & Contents \\
\midrule
Target/junction & 7 & centroids, areas, detection flag \\
Depth geometry & 6 & center/mean/max, left/center/right \\
Motor position & 3 & $M_0,M_1,M_2$ angles \\
Motor velocity & 3 & finite-difference angular velocity \\
Timing & 1 & frame interval $\Delta t$ \\
\midrule
Total & 20 & normalized structured state \\
\bottomrule
\end{tabular}
\end{table}

\section{MOSAIC Goal-Conditioned Recurrent Policy}

\subsection{Architecture}

The policy receives an eight-frame window
$\mathbf O_t=[\obs_{t-7},\ldots,\obs_t]$ and a goal
$\goal\in\{0,\ldots,9\}$. Goals comprise exploration, trachea, and eight named
bronchial destinations. A shared multilayer perceptron encodes each observation
through $20\rightarrow128\rightarrow128$ layers with LayerNorm and ELU. A
learned lookup table maps the goal to a 16-dimensional vector, which is
replicated across time and concatenated with every encoded observation.

The $144$-dimensional sequence is processed by a two-layer GRU with hidden size
256 and inter-layer dropout 0.1:
\begin{align}
 \mathbf x_t &= [f_{\rm obs}(\tilde\obs_t); f_{\rm goal}(\goal)],\\
 \hid_t &= \mathrm{GRU}(\mathbf x_t,\hid_{t-1}).
\end{align}
The final state feeds an action head
$256\rightarrow128\rightarrow3$ with $\tanh$ output and a task head
$256\rightarrow64\rightarrow2$. The action head predicts normalized increments
for insertion and two bending axes. The task head distinguishes navigation from
obstruction intervention and acts as an auxiliary representation constraint.

\subsection{Learning Objective and Data Split}

The action and auxiliary objectives are
\begin{equation}
 \mathcal L = \|\hat\act_t-\act_t^E\|_2^2+
 \lambda_s\,\mathrm{CE}(\hat{\mathbf z}_t,s_t^E),
\end{equation}
where $\act_t^E$ is the expert motor increment and $s_t^E$ is the task label.
Training uses AdamW, OneCycleLR, mixed precision, gradient clipping, and batch
size 64. Demonstrations are split by whole trajectory (90/10), preventing
adjacent windows from the same run from entering both training and validation.

The repository contains 87 real-robot demonstrations and 30,730 synchronized
frames across seven recorded anatomical targets. Windows are sampled with
$T=8$. Actions remain normalized during learning; deployment scales
$\tanh$ outputs by $[5,3,3]^\circ$ per frame for $M_0/M_1/M_2$ and applies an
additional configurable safety factor.
The normalized action is converted to a safe absolute target as
\begin{equation}
 \mathbf q_{t+1}^{\rm cmd}=\operatorname{clip}
 (\mathbf q_t+\beta\,\boldsymbol\alpha\odot\hat\act_t,
  \mathbf q_{\min},\mathbf q_{\max}),
\end{equation}
where $\boldsymbol\alpha=[5,3,3]^\circ$ and $\beta\in(0,1]$ is the deployment
safety factor.

\subsection{Stateful Real-Time Inference}

The runner maintains a fixed-length deque, the goal token, and recurrent hidden
state. At each control step it (i) obtains the newest synchronized feature
dictionary, (ii) estimates velocities from cached angles and $\Delta t$, (iii)
normalizes the 20-dimensional observation, (iv) performs inference without
gradients, and (v) converts bounded increments to absolute motor targets. When
the goal changes or autonomy restarts, the buffer and hidden state are reset.
Missing fields fall back to neutral values, while the surrounding controller
uses timestamps and consecutive-frame confirmation to prevent stale visual
events from triggering mode changes.

Table~\ref{tab:comparison} places the method among its closest autonomous
bronchoscopy counterparts. It is a design comparison, not a performance
ranking, because the systems use different robots, targets, and test domains.

\begin{table}[t]
\caption{Design Context in Autonomous Bronchoscopy}
\label{tab:comparison}
\centering\small
\resizebox{\columnwidth}{!}{%
\begin{tabular}{lllll}
\toprule
Method & Training & Online input & Memory & Validation \\
\midrule
AI co-pilot \cite{zhang2024aicopilot} & sim. imitation & RGB + command & DAgger & phantom/in vivo \\
BronchoCopilot \cite{zhao2024bronchocopilot} & sim. RL & RGB + pose & latent state & simulation \\
Long--short \cite{wu2026longshort} & hierarchical & RGB + CT goal & dual agent & phantom/ex vivo/in vivo \\
MOSAIC-Broncho & real BC & 20-D + goal & 8-frame GRU & silicone phantom \\
\bottomrule
\end{tabular}}
\end{table}

\section{Experimental Protocol}
\label{sec:protocol}

\subsection{Phantom Evaluation}

The final study will use paired, randomized trials from a fixed start jig on
the same airway phantom. Targets include LMB, RMB, LUB, LLB, RUL, RML, and RLL.
Initial pose, illumination, lens contamination, and mucus placement are varied
according to a predefined schedule. The primary comparison is the proposed
goal-conditioned GRU against (B0) expert teleoperation, (B1) frame-wise behavior
cloning, (B2) GRU without goal conditioning, and (B3) reactive visual servoing.

Primary outcomes are route success and wrong-branch rate. Secondary outcomes
are normalized image-plane center error, wall-contact count and duration,
completion time, action jerk, and inference/control latency. Each target and
condition will report trial count, Wilson 95\% confidence intervals for binary
outcomes, median and interquartile range for skewed continuous metrics, and a
paired test with correction for multiple comparisons. Failures are classified
as wrong branch, timeout, perception loss, contact stop, or operator takeover.

\subsection{Ablations}

Four ablations isolate the source of performance: remove the goal embedding;
replace the GRU with a frame-wise MLP; remove depth/path features; and remove
motor velocity and $\Delta t$. A fifth experiment compares the dual-head loss
against action-only behavior cloning. These ablations test whether improvements
arise from goal conditioning, temporal memory, geometric perception, dynamic
state, or auxiliary supervision rather than parameter count alone.

\section{Discussion}

MOSAIC-Broncho occupies a middle ground between reactive servoing and large
hierarchical agents. It retains explicit geometry and robot state, but learns
how their recent history maps to action. This choice reduces model size and
makes failures easier to trace. It also limits semantic breadth. The policy can
only pursue configured anatomical goals, and its competence is bounded by the
expert trajectories it has seen.

The comparison with recent systems is therefore nuanced. AI co-pilot has
stronger sim-to-real training and in-vivo evidence \cite{zhang2024aicopilot}.
BronchoCopilot demonstrates deeper simulated navigation with multimodal RL
\cite{zhao2024bronchocopilot}. Long--short agents address long-horizon planning
and report ex-vivo and in-vivo experiments \cite{wu2026longshort}. Our current
strength is different: the policy learns continuous three-axis control from
real robot demonstrations and operates on a compact, inspectable temporal
state. This is useful, but it is not yet evidence of state-of-the-art navigation
performance.

The present evidence has limitations. Validation MSE measures agreement with
recorded actions, not recovery from covariate shift. Seven of ten configured
goals have demonstrations, class balance is uneven, and perfect auxiliary
accuracy may reflect label simplicity. Monocular depth can drift under fluid,
specularity, and deformation. Finally, physical route success, contacts, and
latency distributions remain to be inserted from randomized trial logs before
submission. These limitations are made explicit to prevent software tests or
qualitative videos from being misreported as efficacy evidence.

\section{Conclusion}

We presented MOSAIC-Broncho, a compact imitation-control framework for
goal-directed robotic bronchoscopy. Its key idea is straightforward: make the
navigation state explicit, then reason over time. A 20-D observation captures
monocular airway cues and robot dynamics. An eight-frame, goal-conditioned GRU
turns that history into bounded three-axis commands. The current artifacts
establish the data scale, model design, and held-out prediction error. The
locked phantom protocol defines the remaining evidence needed for submission.
If the physical trials confirm robust target arrival, the system will show that
useful bronchoscopic autonomy does not always require a large end-to-end model.

\section*{Acknowledgment}
For double-anonymous review, institutional support is omitted. OpenAI Codex was
used to assist with manuscript organization, language drafting, and LaTeX
preparation. All technical claims, citations, equations, and reported values
were checked by the authors against the project code and saved artifacts.

\bibliographystyle{IEEEtran}
\bibliography{refs}

\end{document}
